\documentclass[11pt]{article}

\usepackage[a4paper,margin=1in]{geometry}
\usepackage{amsmath,amssymb}
\usepackage{enumitem}
\usepackage{fancyhdr}
\usepackage{booktabs} % For better looking tables
\usepackage{xcolor} 

\pagestyle{fancy}
\fancyhf{}
\lhead{CSE 400: Fundamentals of Probability in Computing}
\rhead{Lecture 3 Scribe}
\cfoot{\thepage}

\begin{document}

\begin{center}
    {\Large \textbf{Lecture Scribe}} \\
    \vspace{0.3cm}
    \begin{tabular}{rl}
        \textbf{Course:} & CSE400 – Fundamentals of Probability in Computing \\
        \textbf{Lecture Number:} & Lecture 3 \\
        \textbf{Lecture Topic:} & Introduction to Probability Theory \\
        \textbf{Instructor:} & Dr. Dhaval Patel \\
        \textbf{Date:} & January 13, 2026 \\
        \textbf{Audience:} & 2nd Year B.Tech (CSE) \\
        \textbf{Source:} & Lecture 3 Slides (L3.pdf)
    \end{tabular}
\end{center}
\hrule
\vspace{0.5cm}


\section{Opening Context (as per lecture motivation)}
The lecture formally introduces probability theory as a mathematical framework to reason about uncertainty, which naturally arises in:
\begin{itemize}
    \item Daily-life decision making
    \item Engineering systems
    \item Computing systems where outcomes are not deterministic
\end{itemize}
The emphasis in this lecture is on:
\begin{itemize}
    \item Formal definitions
    \item Mathematical structure
    \item Foundations required for later topics (random variables, distributions, algorithms)
\end{itemize}

\section{Random Experiments}

\subsection{Definition}
\textbf{Definition (as stated in lecture):} \\
A random experiment is a process or procedure which, when performed under identical conditions, results in one outcome from a set of possible outcomes, where the exact outcome cannot be predicted with certainty in advance.

\subsection{Key Characteristics}
\begin{itemize}
    \item Repeatable under identical conditions
    \item Outcome is uncertain
    \item Set of all possible outcomes is known
\end{itemize}

\subsection{Examples (from lecture discussion)}
\begin{itemize}
    \item Tossing a coin
    \item Rolling a die
    \item Sending a packet over a noisy communication channel
    \item Sensor measurement under noise
\end{itemize}

\section{Sample Space}

\subsection{Definition}
\textbf{Definition:} \\
The sample space, denoted by $\Omega$, is the set of all possible outcomes of a random experiment.
\[ \Omega = \{ \text{all possible outcomes} \} \]

\subsection{Examples}
\begin{itemize}
    \item \textbf{Single coin toss:} $\Omega = \{H, T\}$
    \item \textbf{Rolling a six-sided die:} $\Omega = \{1, 2, 3, 4, 5, 6\}$
    \item \textbf{Two coin tosses:} $\Omega = \{HH, HT, TH, TT\}$
\end{itemize}

\section{Events}

\subsection{Definition}
\textbf{Definition:} \\
An event is any subset of the sample space $\Omega$.
\[ A \subseteq \Omega \]

\subsection{Types of Events}
\begin{itemize}
    \item \textbf{Simple (Elementary) Event:} Contains exactly one outcome
    \item \textbf{Compound Event:} Contains more than one outcome
\end{itemize}

\subsection{Examples}
Let $\Omega = \{1, 2, 3, 4, 5, 6\}$.
\begin{itemize}
    \item Event A = ``Even number occurs'' $\implies A = \{2, 4, 6\}$
    \item Event B = ``Number greater than 4'' $\implies B = \{5, 6\}$
\end{itemize}

\section{Set Operations on Events}
Events follow set-theoretic operations, which are used heavily later in probability derivations.

\subsection{Union}
\[ A \cup B = \{ \omega : \omega \in A \text{ or } \omega \in B \} \]

\subsection{Intersection}
\[ A \cap B = \{ \omega : \omega \in A \text{ and } \omega \in B \} \]

\subsection{Complement}
\[ A^c = \Omega \setminus A \]

\section{Probability Measure}

\subsection{Definition}
\textbf{Definition (informal, as introduced):} \\
Probability assigns a numerical value to events indicating how likely they are to occur. \\
Formally, probability is a function:
\[ P : F \to [0, 1] \]
where:
\begin{itemize}
    \item $F$ is a collection of events
    \item $P(A)$ denotes probability of event A
\end{itemize}

\section{Axioms of Probability (Kolmogorov Axioms)}
These axioms form the foundation of all probability theory and were explicitly emphasized as non-negotiable assumptions.

\subsection{Axiom 1: Non-negativity}
For any event A:
\[ P(A) \ge 0 \]

\subsection{Axiom 2: Normalization}
\[ P(\Omega) = 1 \]
\textbf{Interpretation:} Some outcome in $\Omega$ must occur.

\subsection{Axiom 3: Additivity}
For mutually exclusive events $A_1, A_2, \dots$:
\[ P\left(\bigcup_i A_i\right) = \sum_i P(A_i) \]
\textbf{Condition:} $A_i \cap A_j = \emptyset$ for $i \neq j$.

\section{Immediate Consequences of Axioms}

\subsection{Probability of Empty Set}
\[ P(\emptyset) = 0 \]
\textbf{Reasoning:}
\[ \Omega = \Omega \cup \emptyset \implies P(\Omega) = P(\Omega) + P(\emptyset) \implies P(\emptyset) = 0 \]

\subsection{Complement Rule}
\[ P(A^c) = 1 - P(A) \]
\textbf{Derivation:}
\[ A \cup A^c = \Omega, \quad A \cap A^c = \emptyset \]
Using additivity:
\[ P(A) + P(A^c) = P(\Omega) = 1 \]

\section{Finite Sample Space with Equally Likely Outcomes}

\subsection{Assumption}
\begin{itemize}
    \item Sample space $\Omega$ is finite
    \item All outcomes are equally likely
\end{itemize}

\subsection{Probability Formula}
For event A:
\[ P(A) = \frac{|A|}{|\Omega|} \]
where:
\begin{itemize}
    \item $|A|$ = number of outcomes in A
    \item $|\Omega|$ = total number of outcomes
\end{itemize}

\section{Worked Example (from lecture style)}
\textbf{Problem:} \\
A fair die is rolled once. Find the probability of getting a number greater than 3.

\textbf{Step-by-Step Solution:} \\
\textbf{Step 1: Define sample space}
\[ \Omega = \{1, 2, 3, 4, 5, 6\} \]
\textbf{Step 2: Define event}
\[ A = \{4, 5, 6\} \]
\textbf{Step 3: Apply formula}
\[ P(A) = \frac{|A|}{|\Omega|} = \frac{3}{6} = \frac{1}{2} \]

\section{Logical Dependencies (Important for Exams)}
\begin{enumerate}
    \item Random experiment $\to$ Sample space
    \item Sample space $\to$ Events
    \item Events + axioms $\to$ Probability rules
    \item All later topics rely on these definitions
\end{enumerate}

\section{What This Lecture Sets Up}
This lecture is the foundation for:
\begin{itemize}
    \item Conditional probability
    \item Random variables
    \item PMF / PDF
    \item Bayes' theorem
    \item Randomized algorithms
\end{itemize}
Without clear understanding of $\Omega$, events, and axioms, later material becomes mechanical and error-prone.

\section{Summary \& Key Takeaways (as emphasized)}
\begin{itemize}
    \item Probability is a mathematical model, not guesswork
    \item Axioms define what probability must satisfy
    \item Events are sets, and probability obeys set logic
    \item Finite, equally-likely model is a special case, not general
\end{itemize}

\newpage
\section*{APPENDIX: Quick Reference (Exam-Ready)}

\subsection*{Core Definitions}
\begin{table}[h!]
    \centering
    \renewcommand{\arraystretch}{1.5}
    \begin{tabular}{|l|l|}
        \hline
        \textbf{Concept} & \textbf{Definition} \\ \hline
        Random Experiment & Process with uncertain outcome \\ \hline
        Sample Space ($\Omega$) & Set of all possible outcomes \\ \hline
        Event & Subset of $\Omega$ \\ \hline
        Probability & Function from events to $[0,1]$ \\ \hline
    \end{tabular}
\end{table}

\subsection*{Axioms}
\begin{itemize}
    \item $P(A) \ge 0$
    \item $P(\Omega) = 1$
    \item Additivity for mutually exclusive events
\end{itemize}

\subsection*{Key Formulas}
\begin{itemize}
    \item $P(A^c) = 1 - P(A)$
    \item $P(A) = |A|/|\Omega|$ (finite, equally likely)
\end{itemize}

\subsection*{Common Exam Warnings}
\begin{itemize}
    \item \textbf{[Warning]} Probability axioms are assumptions, not derived
    \item \textbf{[Warning]} Events are sets, not outcomes
    \item \textbf{[Warning]} Equally-likely assumption must be explicitly valid
\end{itemize}

\end{document}